%========================================
% MatinBook - Stage 08: Code System Test
% Test: Code Listings + Inline Code + Captions + References
% Purpose: Validate minted integration and code display features
%========================================

% Add parent directory to search path
\makeatletter
\def\input@path{{../}}
\makeatother

\documentclass{matinbook}

\title{آزمون سیستم نمایش کد}
\author{تیم توسعه MatinBook}
\date{تابستان ۱۴۰۵}

\begin{document}
    
    %----------------------------------------
    % Title Page
    %----------------------------------------
    \maketitle
    
    %----------------------------------------
    % Chapter 1: Python Programming
    %----------------------------------------
    \chapter{برنامه‌نویسی پایتون}
    
    \section{مبانی زبان}
    
    پایتون\index{پایتون} یک زبان برنامه‌نویسی سطح بالا،
    تفسیری و همه‌منظوره است که بر خوانایی کد تأکید دارد.
    در ادامه، مفاهیم پایه را با کد بررسی می‌کنیم.
    
    \subsection{متغیرها و انواع داده}
    
    در \coderef{code:variables} با انواع داده پایه در پایتون
    آشنا می‌شویم:
    
    \begin{latin}
        \begin{minted}{python}
            # Basic data types in Python
            age: int = 25                          # Integer type
            temperature: float = 36.5              # Floating-point type
            name: str = "Ali"                      # String type
            is_active: bool = True                 # Boolean type
            
            # Compound data structures
            scores: list = [18, 17, 20, 19]       # List - ordered, mutable
            coordinates: tuple = (35.7, 51.4)      # Tuple - ordered, immutable
            student: dict = {                       # Dictionary - key-value pairs
                "name": "Ali",
                "age": 20,
                "grades": [18, 17, 20]
            }
            
            # Set - unordered collection of unique elements
            unique_ids: set = {1001, 1002, 1003}
        \end{minted}
    \end{latin}
    \codecaption{انواع داده در پایتون}
    \label{code:variables}
    
    \subsection{ساختارهای شرطی}
    
    استفاده از دستور \inlcode{if-elif-else} در پایتون:
    
    \begin{latin}
        \begin{minted}{python}
            def get_grade(score: float) -> str:
            """
            Convert numerical score to letter grade.
            
            Args:
            score: Numerical score between 0 and 100.
            
            Returns:
            Letter grade (A, B, C, D, or F).
            
            Examples:
            >>> get_grade(85)
            'B'
            >>> get_grade(92)
            'A'
            """
            if score >= 90:
            return 'A'
            elif score >= 80:
            return 'B'
            elif score >= 70:
            return 'C'
            elif score >= 60:
            return 'D'
            else:
            return 'F'
            
            
            # Create a list of students with their scores
            students = [
            {"name": "Ali", "score": 85},
            {"name": "Sara", "score": 92},
            {"name": "Reza", "score": 78},
            ]
            
            # Calculate and display grades for each student
            for student in students:
            grade = get_grade(student["score"])
            print(f"{student['name']}: {student['score']} -> {grade}")
        \end{minted}
    \end{latin}
    \codecaption{ساختار شرطی و حلقه در پایتون}
    \label{code:conditions}
    
    \subsection{توابع و بازگشت}
    
    \begin{latin}
        \begin{minted}{python}
            def factorial(n: int) -> int:
            """
            Calculate the factorial of n using recursion.
            
            Args:
            n: A non-negative integer.
            
            Returns:
            The factorial of n (n!).
            
            Examples:
            >>> factorial(5)
            120
            >>> factorial(0)
            1
            """
            # Base case: 0! = 1 and 1! = 1
            if n <= 1:
            return 1
            # Recursive case: n! = n * (n-1)!
            return n * factorial(n - 1)
            
            
            def fibonacci(n: int) -> int:
            """
            Calculate the nth Fibonacci number using recursion.
            
            Args:
            n: The index (0-based) in the Fibonacci sequence.
            
            Returns:
            The nth Fibonacci number.
            
            Examples:
            >>> fibonacci(0)
            0
            >>> fibonacci(10)
            55
            """
            # Base cases
            if n <= 0:
            return 0
            elif n == 1:
            return 1
            # Recursive case: F(n) = F(n-1) + F(n-2)
            return fibonacci(n - 1) + fibonacci(n - 2)
            
            
            # Display the first 10 Fibonacci numbers
            for i in range(10):
            print(f"F({i}) = {fibonacci(i)}")
        \end{minted}
    \end{latin}
    \codecaption{توابع بازگشتی در پایتون}
    \label{code:recursion}
    
    همانطور که در \coderef{code:recursion} می‌بینیم،
    پیچیدگی زمانی تابع فیبوناچی بازگشتی $O(2^{n})$ است
    که برای $n$های بزرگ کارآمد نیست.
    
    \subsection{برنامه‌نویسی شیءگرا}
    
    \begin{latin}
        \begin{minted}{python}
            from dataclasses import dataclass
            from typing import List, Optional
            
            
            @dataclass
            class Student:
            """
            A class representing a university student.
            
            Attributes:
            name: Full name of the student.
            student_id: Unique student identification number.
            grades: List of numerical grades (0-20 scale).
            """
            name: str
            student_id: int
            grades: List[float]
            
            @property
            def average(self) -> float:
            """
            Calculate the weighted average of all grades.
            
            Returns:
            The arithmetic mean of grades, or 0.0 if no grades exist.
            """
            if not self.grades:
            return 0.0
            return sum(self.grades) / len(self.grades)
            
            @property
            def letter_grade(self) -> str:
            """Convert numerical average to letter grade."""
            avg = self.average
            if avg >= 18:
            return 'A'
            elif avg >= 15:
            return 'B'
            elif avg >= 12:
            return 'C'
            elif avg >= 10:
            return 'D'
            else:
            return 'F'
            
            def add_grade(self, grade: float) -> None:
            """
            Add a new grade to the student's record.
            
            Args:
            grade: A numerical grade between 0 and 20.
            
            Raises:
            ValueError: If grade is outside the valid range.
            """
            if 0 <= grade <= 20:
            self.grades.append(grade)
            else:
            raise ValueError(
            f"Grade must be between 0 and 20, got {grade}"
            )
            
            
            # Create student instances and add grades
            student = Student("Ali Alizadeh", 1001, [18, 17, 19])
            student.add_grade(20)
            
            # Display student information
            print(f"Student: {student.name}")
            print(f"ID: {student.student_id}")
            print(f"Grades: {student.grades}")
            print(f"Average: {student.average:.2f}")
            print(f"Letter Grade: {student.letter_grade}")
        \end{minted}
    \end{latin}
    \codecaption{کلاس Student با DataClass و Property}
    \label{code:oop}
    
    \section{کد درون خطی}
    
    می‌توانیم در متن به توابع و متغیرها اشاره کنیم:
    تابع \inlcode{print()} برای چاپ خروجی،
    تابع \inlcode{len()} برای محاسبه طول،
    و \inlcode{range()} برای تولید دنباله اعداد.
    
    همچنین می‌توان به ماژول‌های استاندارد مانند
    \inlcode{collections}، \inlcode{itertools}
    و \inlcode{functools} اشاره کرد.
    
    %----------------------------------------
    % Chapter 2: C++ Programming
    %----------------------------------------
    \chapter{برنامه‌نویسی \lr{C++}}
    
    \section{مفاهیم پیشرفته}
    
    \subsection{اشاره‌گرهای هوشمند}
    
    \begin{latin}
        \begin{minted}{c++}
            #include <iostream>
            #include <memory>
            #include <string>
            #include <vector>
            
            /**
            * A simple resource class that logs creation and destruction.
            * Demonstrates RAII (Resource Acquisition Is Initialization).
            */
            class Resource {
                private:
                std::string name;
                
                public:
                // Constructor - acquires the resource
                explicit Resource(const std::string& n) : name(n) {
                    std::cout << "Resource '" << name << "' created\n";
                }
                
                // Destructor - releases the resource
                ~Resource() {
                    std::cout << "Resource '" << name << "' destroyed\n";
                }
                
                // Use the resource
                void use() const {
                    std::cout << "Using resource: " << name << "\n";
                }
                
                // Getter for the resource name
                std::string getName() const {
                    return name;
                }
            };
            
            int main() {
                std::cout << "=== Unique Pointer Demo ===\n";
                
                // unique_ptr: exclusive ownership, cannot be copied
                auto dbConnection = std::make_unique<Resource>("Database-Connection");
                dbConnection->use();
                
                // Transfer ownership using std::move
                auto dbConnection2 = std::move(dbConnection);
                // dbConnection is now nullptr
                if (!dbConnection) {
                    std::cout << "Original pointer is now null\n";
                }
                dbConnection2->use();
                
                std::cout << "\n=== Shared Pointer Demo ===\n";
                
                // shared_ptr: shared ownership with reference counting
                auto cache = std::make_shared<Resource>("Cache-Service");
                std::cout << "Reference count: " << cache.use_count() << "\n";
                
                {
                    // Create another shared_ptr pointing to the same resource
                    auto cacheCopy = cache;
                    std::cout << "Reference count: " << cache.use_count() << "\n";
                    cacheCopy->use();
                }  // cacheCopy goes out of scope, reference count decreases
                
                std::cout << "Reference count: " << cache.use_count() << "\n";
                cache->use();
                
                std::cout << "\n=== End of main ===\n";
                return 0;
            }  // All smart pointers are destroyed here
        \end{minted}
    \end{latin}
    \codecaption{اشاره‌گرهای هوشمند و \lr{RAII} در \lr{C++}}
    \label{code:smart-pointers}
    
    \subsection{الگوها (\lr{Templates})}
    
    \begin{latin}
        \begin{minted}{c++}
            #include <iostream>
            #include <vector>
            #include <stdexcept>
            #include <string>
            #include <concepts>
            
            /**
            * Concept: a type must be comparable using < and > operators.
            * This ensures template type safety at compile time.
            */
            template<typename T>
            concept Comparable = requires(T a, T b) {
                { a < b } -> std::convertible_to<bool>;
                { a > b } -> std::convertible_to<bool>;
                { a == b } -> std::convertible_to<bool>;
            };
            
            /**
            * A generic Stack data structure implemented using std::vector.
            * 
            * @tparam T The type of elements stored in the stack.
            *           Must satisfy the Comparable concept.
            */
            template<typename T>
            requires Comparable<T>
            class Stack {
                private:
                std::vector<T> data;  // Underlying container
                
                public:
                /**
                * Push a new element onto the top of the stack.
                * 
                * @param value The element to push (copied).
                */
                void push(const T& value) {
                    data.push_back(value);
                }
                
                /**
                * Remove and return the top element of the stack.
                * 
                * @return The top element.
                * @throws std::runtime_error if the stack is empty.
                */
                T pop() {
                    if (data.empty()) {
                        throw std::runtime_error("Stack underflow: cannot pop from empty stack");
                    }
                    T value = data.back();
                    data.pop_back();
                    return value;
                }
                
                /**
                * Return a reference to the top element without removing it.
                * 
                * @return Const reference to the top element.
                * @throws std::runtime_error if the stack is empty.
                */
                const T& top() const {
                    if (data.empty()) {
                        throw std::runtime_error("Stack is empty: no top element");
                    }
                    return data.back();
                }
                
                // Capacity queries
                bool empty() const { return data.empty(); }
                size_t size() const { return data.size(); }
                
                /**
                * Remove all elements from the stack.
                */
                void clear() {
                    data.clear();
                }
            };
            
            int main() {
                // Create a stack of integers
                Stack<int> intStack;
                
                // Push elements onto the stack
                intStack.push(10);
                intStack.push(20);
                intStack.push(30);
                
                std::cout << "Stack size: " << intStack.size() << "\n";
                
                // Pop and display all elements (LIFO order)
                std::cout << "Popping elements: ";
                while (!intStack.empty()) {
                    std::cout << intStack.pop() << " ";
                }
                std::cout << "\n";  // Output: 30 20 10
                
                // Create a stack of strings
                Stack<std::string> strStack;
                strStack.push("Hello");
                strStack.push("World");
                
                std::cout << "Top string: " << strStack.top() << "\n";
                
                return 0;
            }
        \end{minted}
    \end{latin}
    \codecaption{الگوها (\lr{Templates}) و مفاهیم (\lr{Concepts}) در \lr{C++20}}
    \label{code:templates}
    
    %----------------------------------------
    % Chapter 3: Algorithm Implementation
    %----------------------------------------
    \chapter{پیاده‌سازی الگوریتم‌ها}
    
    \section{الگوریتم‌های جستجو}
    
    \subsection{جستجوی دودویی}
    
    \begin{latin}
        \begin{minted}{python}
            from typing import List, Optional
            
            
            def binary_search(arr: List[int], target: int) -> Optional[int]:
            """
            Perform binary search on a sorted array.
            
            Binary search repeatedly divides the search interval in half.
            If the target value is less than the middle element, search
            the left half; otherwise, search the right half.
            
            Args:
            arr: A sorted list of integers (ascending order).
            target: The value to search for.
            
            Returns:
            The index of the target if found, None otherwise.
            
            Time Complexity: O(log n)
            Space Complexity: O(1)
            
            Examples:
            >>> binary_search([1, 3, 5, 7, 9], 5)
            2
            >>> binary_search([1, 3, 5, 7, 9], 6)
            None
            """
            # Initialize the search boundaries
            left, right = 0, len(arr) - 1
            
            # Continue searching while the interval is valid
            while left <= right:
            # Calculate the middle index (avoids overflow)
            mid = left + (right - left) // 2
            
            # Check if target is present at mid
            if arr[mid] == target:
            return mid
            # If target is greater, ignore the left half
            elif arr[mid] < target:
            left = mid + 1
            # If target is smaller, ignore the right half
            else:
            right = mid - 1
            
            # Target was not found in the array
            return None
            
            
            # Test the binary search implementation
            if __name__ == "__main__":
            sorted_array = [1, 3, 5, 7, 9, 11, 13, 15, 17, 19]
            
            # Test cases
            test_cases = [7, 1, 19, 8]
            
            for target in test_cases:
            result = binary_search(sorted_array, target)
            if result is not None:
            print(f"Target {target} found at index {result}")
            else:
            print(f"Target {target} not found")
        \end{minted}
    \end{latin}
    \codecaption{پیاده‌سازی جستجوی دودویی با مستندات کامل}
    \label{code:binary-search}
    
    \subsection{مرتب‌سازی سریع}
    
    \begin{latin}
        \begin{minted}{python}
            from typing import List
            
            
            def quicksort(arr: List[int]) -> List[int]:
            """
            Sort a list using the quicksort algorithm.
            
            Quicksort is a divide-and-conquer algorithm that works by:
            1. Selecting a pivot element from the array.
            2. Partitioning the other elements into two sub-arrays
            according to whether they are less than or greater than the pivot.
            3. Recursively sorting the sub-arrays.
            
            Args:
            arr: A list of integers to be sorted.
            
            Returns:
            A new sorted list in ascending order.
            
            Time Complexity: O(n log n) average, O(n²) worst case
            Space Complexity: O(n) for the recursive call stack
            
            Examples:
            >>> quicksort([64, 34, 25, 12, 22, 11, 90])
            [11, 12, 22, 25, 34, 64, 90]
            >>> quicksort([])
            []
            """
            # Base case: arrays with 0 or 1 element are already sorted
            if len(arr) <= 1:
            return arr
            
            # Choose the middle element as the pivot
            # (avoids worst-case for already sorted arrays)
            pivot = arr[len(arr) // 2]
            
            # Partition the array into three parts
            left = [x for x in arr if x < pivot]    # Elements smaller than pivot
            middle = [x for x in arr if x == pivot]  # Elements equal to pivot
            right = [x for x in arr if x > pivot]   # Elements larger than pivot
            
            # Recursively sort left and right, then combine
            return quicksort(left) + middle + quicksort(right)
            
            
            # Demonstrate the quicksort algorithm
            if __name__ == "__main__":
            # Test with an unsorted array
            unsorted = [64, 34, 25, 12, 22, 11, 90, 5]
            print(f"Original array: {unsorted}")
            
            sorted_arr = quicksort(unsorted)
            print(f"Sorted array:   {sorted_arr}")
            
            # Verify the result
            expected = sorted(unsorted)
            assert sorted_arr == expected, "Sorting failed!"
            print("Verification passed: array is correctly sorted")
        \end{minted}
    \end{latin}
    \codecaption{پیاده‌سازی مرتب‌سازی سریع با تأیید صحت}
    \label{code:quicksort}
    
    %----------------------------------------
    % Chapter 4: Code Display Features
    %----------------------------------------
    \chapter{ویژگی‌های نمایش کد}
    
    \section{تنظیمات نمایش}
    
    \lr{MatinBook} از بسته \lr{minted} با تنظیمات زیر
    برای نمایش کد استفاده می‌کند:
    
    \begin{itemize}
        \item \textbf{فونت:} \lr{JetBrains Mono} با اندازه \inlcode{\small}
        \item \textbf{شماره خطوط:} فعال (\texttt{linenos=true})
        \item \textbf{شکستن خطوط:} فعال (\texttt{breaklines=true})
        \item \textbf{کادر:} خط تکی دور کد (\texttt{frame=single})
        \item \textbf{رنگ کادر:} \texttt{matinbordergray}
        \item \textbf{پس‌زمینه:} \texttt{matinlightgray!20}
        \item \textbf{اندازه Tab:} ۴ فاصله (\texttt{tabsize=4})
        \item \textbf{حذف تورفتگی مشترک:} فعال (\texttt{autogobble=true})
    \end{itemize}
    
    \section{ارجاع به کدها}
    
    قابلیت ارجاع به قطعات کد با استفاده از \texttt{\textbackslash coderef}:
    
    \begin{itemize}
        \item \coderef{code:variables} --- انواع داده در پایتون
        \item \coderef{code:conditions} --- ساختارهای شرطی
        \item \coderef{code:recursion} --- توابع بازگشتی
        \item \coderef{code:oop} --- برنامه‌نویسی شیءگرا
        \item \coderef{code:smart-pointers} --- اشاره‌گرهای هوشمند \lr{C++}
        \item \coderef{code:templates} --- الگوها و مفاهیم \lr{C++20}
        \item \coderef{code:binary-search} --- جستجوی دودویی
        \item \coderef{code:quicksort} --- مرتب‌سازی سریع
    \end{itemize}
    
    \section{مقایسه زبان‌ها}
    
    \begin{latin}
        \begin{minted}{text}
            Language     | Typing    | Paradigm            | Speed    | Use Case
            -------------|-----------|---------------------|----------|------------------
            Python       | Dynamic   | Multi-paradigm      | Medium   | Data Science, AI
            C++          | Static    | Multi-paradigm      | Fast     | Systems, Games
            Rust         | Static    | Functional/Imp      | Fast     | Systems, Web
            JavaScript   | Dynamic   | Event-driven        | Medium   | Web Development
            Go           | Static    | Concurrent          | Fast     | Cloud, Network
        \end{minted}
    \end{latin}
    \codecaption{مقایسه زبان‌های برنامه‌نویسی}
    \label{code:comparison}
    
    %----------------------------------------
    % Summary
    %----------------------------------------
    \chapter{خلاصه نتایج}
    
    \section{وضعیت سیستم کد}
    
    \begin{table}[h]
        \centering
        \caption{وضعیت قابلیت‌های نمایش کد}
        \label{tab:code-status}
        \begin{tabular}{|c|C{6cm}|c|}
            \hline
            \textbf{قابلیت} & \textbf{توضیحات} & \textbf{وضعیت} \\
            \hline
            \lr{Syntax Highlighting} & پشتیبانی از ۳۰۰+ زبان برنامه‌نویسی & \textcolor{matingreen}{✅} \\
            شماره خطوط & نمایش شماره کنار هر خط کد & \textcolor{matingreen}{✅} \\
            کادر & خط تکی با رنگ \lr{matinbordergray} & \textcolor{matingreen}{✅} \\
            شکستن خطوط & خطوط طولانی به طور خودکار می‌شکنند & \textcolor{matingreen}{✅} \\
            کد درون خطی & \texttt{\textbackslash inlcode} برای ارجاع در متن & \textcolor{matingreen}{✅} \\
            ارجاع به کد & \texttt{\textbackslash coderef} با شماره خودکار & \textcolor{matingreen}{✅} \\
            \lr{Type Hints} & نمایش راهنمای نوع در \lr{Python} & \textcolor{matingreen}{✅} \\
            \lr{Docstrings} & مستندات کامل توابع و کلاس‌ها & \textcolor{matingreen}{✅} \\
            فونت کد & \lr{JetBrains Mono} با لیگچرهای برنامه‌نویسی & \textcolor{matingreen}{✅} \\
            \hline
        \end{tabular}
    \end{table}
    
    \section{نتیجه}
    
    \textbf{ماژول نمایش کد \lr{MatinBook} نسخه ۱ با موفقیت آزمایش شد.}
    کدهای \lr{Python}، \lr{C++} و \lr{text} با رنگ‌آمیزی مناسب،
    مستندات کامل و قابلیت ارجاع به درستی کار می‌کنند. 🎉
    
\end{document}